\section{Method}
\subsection{Learning the Hidden Units for HuBERT}
An acoustic model trained on text and speech pairs provides pseudo-phonetic labels for each frame via forced alignment in semi-supervised learning. On the contrary, the self-supervised representation learning setup has access to speech-only data. Nevertheless, simple discrete latent variable models such as k-means and Gaussian mixture models (GMMs) infer hidden units that exhibit non-trivial correlation with the underlying acoustic units~\cite{lee2012nonparametric} (see also Table~\ref{tab:loss}). More advanced systems can achieve better acoustic unit discovery performance using better graphical models \cite{ondel2016variational, ebbers2017hidden} or parameterizes the distributions with more powerful neural network models~\cite{hsu2017learning, hsu2017unsupervised, chorowski2019unsupervised, khurana2019factorial, khurana2020convolutional}.

\begin{figure}[th]
  \centering
  \includegraphics[width=\linewidth]{figures/hubert_arch}\caption{The HuBERT approach predicts hidden cluster assignments of the masked frames ($y_2, y_3, y_4$ in the figure) generated by one or more iterations of k-means clustering.}
  \label{fig:arch}
\end{figure}
Inspired by this, we propose to use acoustic unit discovery models to provide frame-level targets. 
Let $X$ denote a speech utterance $X = [x_1, \cdots, x_T]$ of $T$ frames.
Discovered hidden units are denoted with $h(X) = Z = [z_1, \cdots, z_T]$, where $z_t \in [C]$ is a $C$-class categorical variable and $h$ is a clustering model, e.g. k-means. 

\subsection{Representation Learning via Masked Prediction}\label{sec:maskpred}
Let $M \subset [T]$ denote the set of indices to be masked for a length-$T$ sequence $X$, and $\tilde{X} = r(X, M)$ denote a corrupted version of $X$ where $x_t$ is replaced with a mask embedding $\tilde{x}$ if $t \in M$. A masked prediction model $f$ takes as input $\tilde{X}$ and predicts a distribution over the target indeces at each timestep $p_f(\cdot \mid \tilde{X}, t)$. There are two decisions to be made for masked prediction: \textit{how to mask} and \textit{where to apply the prediction loss}. 

Regarding the first decision, we adopt the same strategies used in SpanBERT~\cite{joshi2020spanbert} and wav2vec 2.0~\cite{baevski2020wav2vec} for mask generation, where $p$\% of the timesteps are randomly selected as start indices, and spans of $l$ steps are masked. To address the second decision, we denote the cross-entropy loss computed over masked and unmasked timesteps as $L_m$ and $L_u$, respectively. $L_m$ is defined as:
\begin{equation}
    L_m(f; X, M, Z) = \sum_{t \in M} \log p_f(z_t \mid \tilde{X}, t),
\end{equation}
and $L_u$ is of the same form except that it sums over $t \not\in M$.
The final loss is computed as a weighted sum of the two terms: $L = \alpha L_m + (1-\alpha)L_u$. In the extreme case when $\alpha = 0$, the loss is computed over the unmasked timesteps, which is similar to acoustic modeling in hybrid speech recognition systems \cite{young1996large, abdel2012applying, povey2005discriminative, bourlard2012connectionist}. In our setup, this limits the learning process to mimicking the clustering model. 

In the other extreme with $\alpha=1$, the loss is only computed over the masked timesteps where the model has to predict the targets corresponding to the unseen frames from context, analogous to language modeling. It forces the model to learn both the acoustic representation of unmasked segments and the long-range temporal structure of the speech data. We hypothesize that the setup with $\alpha=1$ is more resilient to the quality of cluster targets, which is demonstrated in our experiments (see Table~\ref{tab:loss}).


\subsection{Learning with Cluster Ensembles}
A simple idea to improve target quality is to utilize multiple clustering models. While an individual clustering model may perform terribly, cluster ensembles can provide complementary information to facilitate representation learning. For example, an ensemble of k-means models with different codebook sizes can create targets of different granularity, from manner classes (vowel/consonant) to sub-phone states (senones). 
To extend the proposed framework, let $Z^{(k)}$ be the target sequences generated by the $k$-th clustering model. We can now re-write $L_m$ as:
\begin{equation}
    L_m(f; X, \{ Z^{(k)} \}_k, M) = 
    \sum_{t \in M} \sum_{k} \log p_f^{(k)}(z_t^{(k)} \mid \tilde{X}, t), \label{eq:our_obj}
\end{equation}
and similarly for the unmasked loss $L_u$. This is analogous to multi-task learning, but with tasks created by unsupervised clustering.

Additionally, ensembling is intriguing because it can be used alongside product quantization (PQ)~\cite{gray1998quantization}, where a feature space is partitioned into multiple subspaces, and each subspace is quantized separately. PQ allows effective Euclidean distance-based quantization such as k-means for high-dimensional features and heterogeneous features whose scale differs significantly between subspaces. In this case, the theoretical size of the target space is the product of all codebooks' sizes.


\subsection{Iterative Refinement of Cluster Assignments}
In addition to using cluster ensembles, another direction for improved representation is \textit{refining} the cluster assignments throughout the learning process. Since we expect a pre-trained model to provide better representations than the raw acoustic feature such as MFCCs, we can create a new generation of clusters by training a discrete latent model over the learned latent representations. The learning process then proceeds with the newly discovered units.


\subsection{Implementation}\label{sec:impl}
Our pre-trained models follows the wav2vec 2.0 architecture~\cite{baevski2020wav2vec}, with a convolutional waveform encoder, a BERT encoder~\cite{devlin2018bert}, a projection layer and a code embedding layer. We consider HuBERT in three different configurations: \textsc{Base}, \textsc{Large}, and \textsc{X-Large}. The fisrt two follow the architectures of wav2vec 2.0 \textsc{Base} and \textsc{Large} closely. The \textsc{X-Large} architecture expands the model size to about 1 billion parameters, similar to the size of the Conformer XXL model in ~\cite{zhang2020pushing}. 
The waveform encoder is identical for all the three configurations, which is composed of seven 512-channel layers with strides [5,2,2,2,2,2,2] and kernel widths [10,3,3,3,3,2,2]. The BERT encoder consists of many identical transformer blocks, whose parameters along with the parameter of the subsequent projection layer are specified in Table~\ref{tab:arch}.

\begin{table}[ht]
    \centering
    \begin{tabular}{cc|ccc}
        \toprule
         & & \textsc{Base} & \textsc{Large} & \textsc{X-Large} \\
        \midrule
        \multirow{3}{*}{CNN Encoder}
        & strides      & \multicolumn{3}{c}{5, 2, 2, 2, 2, 2, 2} \\
        & kernel width & \multicolumn{3}{c}{10, 3, 3, 3, 3, 2, 2} \\
        & channel      & \multicolumn{3}{c}{512} \\
        \midrule
        \multirow{5}{*}{Transformer} 
        & layer        & 12 & 24 & 48 \\
        & embedding dim. & 768 & 1024 & 1280 \\
        & inner FFN dim. & 3072 & 4096 & 5120 \\
        & layerdrop prob  & 0.05 & 0 & 0 \\
        & attention heads & 8 & 16 & 16 \\
        \midrule
        \multirow{1}{*}{Projection} & dim. & 256 & 768 & 1024 \\
        \midrule
        \multicolumn{2}{c|}{Num. of Params} & 95M & 317M & 964M \\
        \bottomrule
        
    \end{tabular}
    \caption{Model architecture summary for \textsc{Base}, \textsc{Large}, and \textsc{X-Large} HuBERT models}
    \label{tab:arch}
\end{table}


The convolutional waveform encoder generates a feature sequence at a 20ms framerate for audio sampled at 16kHz (CNN encoder down-sampling factor is 320x). The audio encoded features are then randomly masked as described in Section \ref{sec:maskpred}. The BERT encoder takes as input the masked sequence and outputs a feature sequence $[o_1, \cdots, o_T]$. The distribution over codewords is parameterized with
\begin{equation}
    p_f^{(k)}(c \mid \tilde{X}, t) = \frac{\exp(\text{sim}(A^{(k)} o_t, e_c) / \tau)} {\sum_{c'=1}^C \exp(\text{sim}(A^{(k)} o_t, e_{c'}) / \tau)},
\end{equation}
where $A$ is the projection matrix, $e_c$ is the embedding for codeword $c$, $\text{sim}(\cdot, \cdot)$ computes the cosine similarity between two vectors, and $\tau$ scales the logit, which is set to 0.1. When cluster ensembles are used, one projection matrix $A^{(k)}$ is applied for each clustering model $k$.

After HuBERT pre-training, We use the connectionist temporal classification (CTC)~\cite{graves2006connectionist} loss for ASR fine-tuning of the whole model weights except the convolutional audio encoder, which remains frozen. The projection layer(s) is removed and replaced with a randomly initialized softmax layer. The CTC target vocabulary includes 26 English characters, a space token, an apostrophe, and a special CTC blank symbol.
